%% Example CV - Darshan Verma
%% Comprehensive overview of all competencies, experience, and achievements.
%% Use as a master reference when tailoring role-specific CVs.
%%
%% SETUP: Run /setup to populate this with your actual information.
%% Compile with: lualatex main_example.tex
%% (pdflatex often fails on modern MiKTeX with fontawesome5 font-expansion errors.)

\documentclass[11pt,a4paper,sans]{moderncv}
\moderncvstyle{banking}
\moderncvcolor{blue}

% Force the name and section headings to render in moderncv blue (color1).
% Default banking leaves them black: moderncvstylebanking.sty's \colorlet
% copies (not aliases) the pre-scheme accent colour, so the name colours are
% frozen before \moderncvcolor runs. Re-let them after. \namefont is the hook
% every name-style macro routes through, so this also works on moderncv 2.3.1
% (Debian/Ubuntu apt), which has no \firstnamestyle/\lastnamestyle at all.
\renewcommand*{\namefont}{\fontsize{34}{36}\bfseries\upshape}
\colorlet{firstnamecolor}{color1}
\colorlet{lastnamecolor}{color1}
\colorlet{namecolor}{color1}
\renewcommand*{\sectionstyle}[1]{{\sectionfont\color{color1}#1}}

\usepackage[utf8]{inputenc}
% pdflatex fallback only (the documented engine is lualatex, which skips this
% branch). Without T1 font encoding pdflatex builds accented letters with
% \accent, and the PDF text layer stores them decomposed - `e` + U+0300 rather
% than U+00E8 - so an ATS keyword match on "Genève" fails while the page looks
% right. moderncv 2.5 loads T1 itself under pdflatex; 2.3.1 (Debian/Ubuntu apt)
% does not. \ifpdftex comes from iftex, which every moderncv version loads.
\ifpdftex\usepackage[T1]{fontenc}\fi
% moderncv loads hyperref itself in an \AtEndPreamble hook, so \hypersetup
% must go in an \AtEndPreamble of our own: on moderncv < 2.4 a top-level
% \usepackage{hyperref} clashes with the class's own
% \RequirePackage[unicode]{hyperref}. From 2.4.0 the class passes its options
% through \PassOptionsToPackage instead, which is what removes that clash.
\AtEndPreamble{\hypersetup{
    colorlinks=true,
    linkcolor=blue,
    filecolor=magenta,
    urlcolor=blue,
    pdftitle={Darshan Verma - CV},
    % Keep pdfpagemode=UseNone: this block runs after moderncv's own
    % \AtEndPreamble (moderncv.cls sets pdfpagemode there), so a FullScreen
    % value here would win and open every CV in fullscreen presentation mode.
    pdfpagemode=UseNone,
}}
\usepackage[scale=0.80]{geometry}
\usepackage{import}
\usepackage{needspace}

% personal data
\name{Darshan}{Verma}
\address{Hyderabad, Telangana, India}{}{}
\phone[mobile]{+91 9925531392}
\email{darshanverma2002@gmail.com}
\extrainfo{\href{https://www.linkedin.com/in/dv64bit}{LinkedIn}, \href{https://darshanvportfolio.vercel.app}{Portfolio}}

\begin{document}

\makecvtitle

% ============================================================
%     PROFILE STATEMENT
% ============================================================

\vspace{6pt}
\small{Product Designer with a Computer Science Engineering background, working on AI-driven enterprise SaaS for clients including Moody's and Marvell Technology. Primary hands-on designer at Cognida AI, owning discovery, information architecture, user flows, UI, and developer handoff, and founder of the company's first design system. Builds React and Tailwind prototypes of own designs so stakeholders test real flows and engineers get an unambiguous handoff.}

% ============================================================
%     CORE COMPETENCIES
% ============================================================

\section{Core Competencies}
\vspace{1pt}
\begin{itemize}

\item \textbf{Product \& UX Design}: User Research, Usability Testing, Heuristic Evaluation, Information Architecture, User Flows, Site Mapping, Customer Journey Mapping, Wireframing, Prototyping

\item \textbf{Design Systems}: Design tokens, component libraries, theming, design documentation, token-audit tooling

\item \textbf{Collaboration \& Delivery}: Stakeholder Collaboration, Cross-functional Collaboration, Design QA, Developer Handoff

\item \textbf{Design Engineering}: React 18/19, TypeScript, Tailwind CSS, Vite, frontend fundamentals

\item \textbf{AI in Workflow}: Claude Code, Figma Make, ChatGPT, v0, Lovable

\item \textbf{Tools}: Figma, FigJam, Miro, Framer, Webflow, Notion, Jira

\end{itemize}

% ============================================================
%     PROFESSIONAL EXPERIENCE
% ============================================================

\section{Professional Experience}
\vspace{3pt}
\begin{itemize}

% --- Most Recent Role ---
\needspace{5\baselineskip}
\item{\cventry{Jun 2025--Present}{UX Designer}{Cognida AI}{Hyderabad, India}{}{\vspace{1pt}
\begin{itemize}
    \item Designed a platform covering the full life cycle of a Moody's credit rating report: design-system primitives, block-based templates, authoring, reviewer comments, approval, and publishing. Replaced a process split between MS Word and email. Specced for 2,000+ users, currently rolling out.
    \item Designed a 25-screen procure-to-pay hub for Marvell Technology, pulling request tracking out of Oracle EBS, Beeline, and Conga into one view. Built the front end in React 18, TypeScript, and Tailwind so stakeholders could walk the real flow, then handed the codebase to engineering for backend integration.
    \item Scoped a 20-capability multi-cloud brief into a 5-capability MVP for InfraWatch, an AWS cost optimization dashboard, built around one verifiable loop. Designed 12 screens, treating loading, empty, error, and insufficient-data states as screens in their own right.
    \item Designed Space Org Console, the system behind InfraWatch: 191 tokens driving dark, light, and print themes from one source, with an audit script that blocks raw hex and off-scale type. Coded the prototype in React 19 and Tailwind v4, including its animated charts.
    \item Established Cognida's first design system: color, typography, and spacing tokens plus a component library built on vetted open-source primitives and restyled to the company's visual language. Used across active product teams.
    \item Owned end-to-end design of CogniVault, an internal knowledge platform. Re-scoped it after stakeholder review, cutting five personas to three plus a config surface and collapsing a 14-tab project view into a compliance-first structure.
\end{itemize}}}

\vspace{3pt}

% --- Previous Role ---
\needspace{5\baselineskip}
\item{\cventry{Oct 2024--May 2025}{Customer Delight Executive}{Zomato}{Hyderabad, India}{}{\vspace{1pt}
\begin{itemize}
    \item Self-funded a UI/UX bootcamp while working full time in customer operations, moving from engineering into design. Resolved customer issues through structured support processes and learned where service design breaks from the complaints side of it.
\end{itemize}}}

\vspace{3pt}

% --- Earlier Role ---
\needspace{5\baselineskip}
\item{\cventry{Jun 2024--Aug 2024}{Product Designer (Intern)}{Wizzer}{Mumbai, India}{}{\vspace{1pt}
\begin{itemize}
    \item Designed the Client Management dashboard for a fintech advisory platform, from user journeys and wireframes to high-fidelity Figma prototypes. Ran usability tests and reworked task flows with QA and developers from the feedback.
\end{itemize}}}

\end{itemize}

% ============================================================
%     EDUCATION
% ============================================================

\section{Education \& Certification}
\vspace{1pt}
\begin{itemize}

\needspace{5\baselineskip}
\item{\cventry{2020--2024}{B.E., Computer Science Engineering}{Ajeenkya DY Patil University}{Pune, India}{}{\vspace{1pt}
CGPA 8.27/10. Specialization in Cloud Technology and Information Security.
}}

\vspace{3pt}

\needspace{5\baselineskip}
\item{\cventry{2025}{UI/UX Design with AI, six-month placement guarantee course}{Internshala Trainings}{Online}{}{\vspace{1pt}
Certificate no. 3vdfkjqfomb -- \href{https://trainings.internshala.com/s/v/3657877/28faa386}{Verify}.
}}

\end{itemize}

% ============================================================
%     LANGUAGES
% ============================================================

\section{Languages}
\vspace{1pt}
\begin{itemize}
\item English (professional working proficiency), Hindi (native), Gujarati (native).
\end{itemize}

% ============================================================
%     REFERENCES
% ============================================================

\section{References}
\vspace{1pt}
\begin{itemize}
\item{Available upon request.}
\end{itemize}

\end{document}
